\section{参考文献与致谢}\label{sec:references}

\subsection{核心参考}

\begin{enumerate}
    \item Anthropic (2026). \textit{Claude Code Source Code Leak — Industry Analysis}. VentureBeat 报道及多篇基于公开制品的架构解读.
    \item Cognition Labs (2024). \textit{Devin: Autonomous Software Engineer}. 长程任务 Agent 范式参考.
    \item Sakana AI (2024). \textit{The AI Scientist}. 自动研究 Agent 的反思与 skill 学习机制.
    \item NousResearch. \textit{Hermes Agent}. \url{https://github.com/NousResearch/hermes-agent}. delegate 子代理隔离与上下文压缩参考.
    \item Park et al. (2023). \textit{Generative Agents: Interactive Simulacra of Human Behavior}. Stanford. SimPredictor 虚拟受众设计参考.
    \item Madaan et al. (2023). \textit{SELF-REFINE: Iterative Refinement with Self-Feedback}. SelfCritiqueLoop 设计参考.
    \item Shinn et al. (2023). \textit{Reflexion: Language Agents with Verbal Reinforcement Learning}. Reflector 模块参考.
    \item Packer et al. (2023). \textit{MemGPT: Towards LLMs as Operating Systems}. MemorySystem 分页参考.
    \item Wang et al. (2023). \textit{Voyager: Open-Ended Embodied Agent}. Skill Library 自进化参考.
    \item Microsoft (2023). \textit{AutoGen: Multi-Agent Conversation Framework}. Forum Debate 设计参考.
    \item LangChain (2024). \textit{LangGraph}. 显式 DAG 编排参考.
    \item crewAIInc. \textit{CrewAI}. 多角色 Crew 协作参考.
    \item Polymarket. Gamma API. \url{https://docs.polymarket.com}. 预测市场数据接入参考.
    \item Kalshi. Trading API. \url{https://docs.kalshi.com}. 合规预测市场对比参考.
\end{enumerate}

\subsection{算法与方法}

\begin{enumerate}
    \setcounter{enumi}{14}
    \item Page (1954). \textit{Continuous Inspection Schemes}. CUSUM 变点检测算法.
    \item Iglewicz \& Hoaglin (1993). \textit{How to Detect and Handle Outliers}. MAD-zscore 鲁棒统计.
    \item Granger (1969). \textit{Investigating Causal Relations by Econometric Models}. Trend Chain 因果链推理.
    \item Platt (1999). \textit{Probabilistic Outputs for Support Vector Machines}. Risk-Reward Calibration.
    \item Kirkpatrick et al. (2017). \textit{Overcoming catastrophic forgetting in neural networks}. Persona Vector EWC 防遗忘.
\end{enumerate}

\subsection{平台与数据源}

\begin{enumerate}
    \setcounter{enumi}{19}
    \item HackerNews API. \url{https://github.com/HackerNews/API}.
    \item 微博 / 抖音 / 百度 / B 站 热搜聚合 API. \url{https://v2.xxapi.cn}.
    \item Manifold Markets API. \url{https://docs.manifold.markets}.
    \item DailyHotApi. \url{https://github.com/imsyy/DailyHotApi} (灵感参考).
\end{enumerate}

\subsection{腾讯生态}

\begin{enumerate}
    \setcounter{enumi}{23}
    \item 腾讯混元大模型. \url{https://hunyuan.tencent.com}. LLM Fallback 之一.
    \item 微信视频号开放平台. \url{https://channels.weixin.qq.com}.
    \item 微信公众平台. \url{https://mp.weixin.qq.com}.
\end{enumerate}

\subsection{灵感来源 (个人项目)}

\begin{enumerate}
    \setcounter{enumi}{26}
    \item Digital Oracle 概念 — 「资本永远先于舆论」哲学.
    \item MiroFish — 群体仿真 Agent 设计.
    \item BettaFish — 多 Agent 辩论. \url{https://github.com/Kocoro-lab/BettaFish}.
\end{enumerate}

\subsection{致谢}

感谢腾讯 PCG 校园 AI 产品创意大赛提供的舞台。感谢 Anthropic 公开的 Claude Code 设计哲学。感谢小米 MiMo 和 MiniMax 提供的 LLM 算力。感谢 Polymarket / HackerNews / Manifold 提供的开放数据。

特别致谢:
\begin{itemize}
    \item 中国科学技术大学计算机学院的开放学习环境
    \item 个人主页 \url{https://bcefghj.github.io/} 维护期间的所有读者反馈
    \item 小红书 \texttt{bcefghj} 账号的粉丝, 你们的真实互动数据是 Cohort Affinity 设计的灵感来源
\end{itemize}

\subsection{开源声明}

Ripple 2.0 在 GitHub 开源 (Apache 2.0). 所有第三方借鉴均在文档中标注来源. 不复制任何泄露代码, 仅借鉴公开文章中的设计思想.

\begin{quote-box}
\textbf{联系作者}\\
姓名: 戴尚好\\
学校: 中国科学技术大学\\
邮箱: \texttt{bcefghj@163.com}\\
个人主页: \url{https://bcefghj.github.io/}\\
GitHub: \url{https://github.com/bcefghj}\\
小红书 ID: \texttt{bcefghj}\\
项目地址: \url{https://github.com/bcefghj/ripple}\\
在线 Demo: \url{http://120.55.247.6/chat}
\end{quote-box}
